\documentclass[a4paper,10pt,twoside]{article}
\usepackage{amsmath}
\usepackage{its_conf}
\begin{document}
\setcounter{section}{0}
\setcounter{figure}{0}
\setcounter{table}{0}
\setcounter{equation}{0}
\setcounter{secnumdepth}{1}
% ======= Начало тезисов ===========

\topic{Нестационарность долгосрочной памяти волатильности финансовых временных рядов и её прогностическая ценность}

\information{Бондаренко~Д.~Д.}
{Кафедра теории вероятностей и математической статистики, Белорусский государственный университет}
{Минск, Республика Беларусь}
{adarka.bon@gmail.com}

\annotation{Рассматривается нестационарность долгосрочной памяти в волатильности финансовых рядов и её связь с механизмами внимания трансформеров. Память волатильности, описываемая моделями FIGARCH, ARFIMA-RV и HAR-RV, меняет силу во времени, тогда как архитектуры Informer, Autoformer, PatchTST используют статичные схемы внимания, не следующие за динамикой памяти. Скользящие оценки ($R/S$, GPH) по десяти активам за 2019--2023~гг. показывают, что параметр $d$ реализованной волатильности значимо нестационарен: CUSUM-тест выявляет сдвиги для 7 из 10 активов, в кризисы память ослабевает (среднее $\hat d$ падает с 0.52 до 0.45). Локальная оценка памяти улучшает HAR-RV по MSE для всех 10 активов ($+7.1\,\%$) и по QLIKE для 8 из 10 ($+1.8\,\%$). Самовнимание интерпретируется как обучаемая автокорреляционная функция; обосновывается адаптивный механизм внимания (Memory-Adaptive Attention) для прогноза волатильности и оценки риска VaR/ES.}

\begin{multicols}{2}

\section*{Введение}

Долгосрочная память в волатильности финансовых рядов является устойчиво наблюдаемым фактом и лежит в основе моделей FIGARCH, ARFIMA-RV и HAR-RV. Однако эмпирические исследования указывают на то, что этот эффект нестационарен: сила памяти, то есть степень влияния далёкого прошлого на текущую динамику, меняется во времени и особенно заметно ослабевает в периоды рыночных шоков.

С точки зрения методов обработки информации задача распадается на два комплементарных пласта --- статистический и нейросетевой. Статистические методы долгой памяти (GPH, $R/S$-анализ, CUSUM) измеряют силу и изменчивость памяти, а нейросетевой аппарат трансформеров --- учитывает её в модели. Механизм самовнимания~[1] строит обучаемую нелинейную автокорреляционную функцию, однако схемы внимания в архитектурах Informer~[2], Autoformer и PatchTST статичны: они задаются априори и не зависят от текущего режима памяти. Если сила памяти меняется во времени, фиксированное внимание рассогласуется со структурой данных и систематически ошибается в моменты смены режима.

Целью работы является проверка гипотезы о нестационарности долгосрочной памяти волатильности, оценка прогностической ценности её локальной оценки и анализ того, как эта нестационарность может быть учтена в механизме внимания трансформера.

\section{Данные и методология}

Эмпирическая база включает три класса активов за 2019--2023~гг.: акции (META, BA, JPM, AAPL, MSFT), криптоактивы (BTC, ETH) и валютные пары (EURUSD, USDJPY, GBPUSD). По дневным ценам закрытия рассчитывались логарифмические доходности $r_t=\ln(P_t/P_{t-1})$ и реализованная волатильность $\mathrm{RV}_t=\sqrt{\sum_{i=0}^{20} r_{t-i}^2}$.

Нестационарность памяти измерялась двумя методами в скользящем окне шириной $W=500$ наблюдений с шагом 25: $R/S$-анализом (показатель Хёрста $H$~[3]) и методом GPH~[4], оценивающим параметр дробной интеграции $d$ регрессией логарифма периодограммы по первым $m=n^{0.5}$ низкочастотным гармоникам.

Структурные сдвиги в $\hat d(t)$ проверялись CUSUM-тестом, различие кризисных и спокойных окон~-- тестом Уэлча. Прогностическая ценность памяти оценивалась сравнением модели HAR-RV~[5] с расширением на признаки $\hat d_t$ и $\Delta \hat d_t$ методом расширяющегося окна; качество измерялось MSE и QLIKE.

Для оценки реакции механизмов внимания на смену режима памяти проведён пилотный эксперимент на минутных котировках: три архитектуры --- полное самовнимание (Standard), разреженное ProbSparse и двухпутевое MultiFeature --- с одинаковой размерностью $d_{\text{model}}=32$, четырьмя головами и двумя слоями обучены на одном наборе признаков; тестовые наблюдения разбивались на режимы высокой и низкой локальной памяти по квартилям $\hat d(t)$.

\section{Результаты}

Долгосрочная память в волатильности проявляется существенно сильнее, чем в доходностях: для реализованной волатильности $\hat d$ лежит в диапазоне $0.42$--$0.56$ ($H\approx0.90$--$0.95$), тогда как для логарифмических доходностей $\hat d\approx0.01$--$0.07$ (таблица~\ref{tab:mem}). Оценки при этом не постоянны --- размах $\hat d(t)$ достигает $0.5$ (для MSFT $0.33$--$0.75$), а на рисунке~\ref{fig:mem} видно, что в цикл повышения ставок 2022~г. $\hat d(t)$ для акций и валютных пар синхронно опускается с $0.65$--$0.75$ до $0.30$--$0.45$: память <<схлопывается>> на фоне кризиса.

\begin{table}[H]
\centering\small
\caption{Локальная память RV: среднее, режимы кризис/спокойствие и значимость сдвига ($W=500$)}
\label{tab:mem}
\begin{tabularx}{\columnwidth}{|X|c|c|c|c|}
\hline
Актив & $\bar{\hat d}$ & $\hat d_{\text{кр.}}$ & $\hat d_{\text{сп.}}$ & $p_{\text{CUSUM}}$ \\
\hline
META   & 0.426 & 0.386 & 0.446 & $<0.001$ \\
BA     & 0.538 & 0.480 & 0.568 & $<0.001$ \\
JPM    & 0.539 & 0.469 & 0.574 & $<0.001$ \\
AAPL   & 0.556 & 0.472 & 0.598 & 0.003 \\
MSFT   & 0.548 & 0.480 & 0.582 & $<0.001$ \\
BTC    & 0.419 & 0.416 & 0.420 & 0.056 \\
ETH    & 0.430 & 0.443 & 0.424 & 0.448 \\
EURUSD & 0.506 & 0.430 & 0.551 & $<0.001$ \\
USDJPY & 0.516 & 0.458 & 0.550 & $<0.001$ \\
GBPUSD & 0.505 & 0.463 & 0.531 & 0.061 \\
\hline
\end{tabularx}
\end{table}

\image{
[width=\columnwidth]{paper_memory.png}
\caption{Скользящие оценки $\hat d(t)$ и $H(t)$ для реализованной волатильности (красным выделены кризисные периоды)\label{fig:mem}}
}

CUSUM-тест обнаруживает структурный сдвиг в $\hat d(t)$ для 7 из 10 активов, а различие средних кризисных и спокойных окон значимо для 7 из 10 (таблица~\ref{tab:mem}): память в кризисные периоды ниже, $\bar{\hat d}$ падает с $0.52$ до $0.45$. Криптоактивы (BTC, ETH) не демонстрируют значимого падения, что согласуется с их особой режимной динамикой.

Добавление признаков $\hat d_t$, $\Delta\hat d_t$ улучшает HAR-RV для всех 10 активов по MSE (средний прирост $+7.1\,\%$) и для 8 из 10 по QLIKE ($+1.8\,\%$) (таблица~\ref{tab:har}); наибольший эффект --- на волатильных инструментах (META $+26.9\,\%$, BTC $+16.1\,\%$).

\begin{table}[H]
\centering\small
\caption{HAR-RV против HAR-RV с локальной памятью}
\label{tab:har}
\begin{tabularx}{\columnwidth}{|X|c|c|}
\hline
Актив & $\Delta$MSE, \% & $\Delta$QLIKE, \% \\
\hline
META   & $+26.9$ & $+3.8$ \\
BA     & $+1.0$  & $+0.6$ \\
JPM    & $+1.0$  & $+4.1$ \\
AAPL   & $+2.2$  & $+0.9$ \\
MSFT   & $+0.3$  & $-5.7$ \\
BTC    & $+16.1$ & $+5.4$ \\
ETH    & $+2.4$  & $+2.1$ \\
EURUSD & $+8.5$  & $+3.8$ \\
USDJPY & $+6.4$  & $+3.1$ \\
GBPUSD & $+5.8$  & $-0.3$ \\
\hline
\end{tabularx}
\end{table}

Пилотный эксперимент на минутных котировках показывает, что статичные механизмы внимания по-разному деградируют при смене режима памяти: разреженное ProbSparse-внимание, будучи лучшим при низкой памяти, при её усилении теряет точность более чем вдвое, тогда как двухпутевое MultiFeature остаётся устойчивее (таблица~\ref{tab:attn}).

\begin{table}[H]
\centering\small
\caption{RMSE механизмов внимания по режимам памяти и его прирост}
\label{tab:attn}
\begin{tabularx}{\columnwidth}{|X|c|c|c|}
\hline
Механизм & низк. & выс. & прирост, \% \\
\hline
Standard & 1.078 & 1.628 & $+51$ \\
ProbSparse & 0.415 & 1.006 & $+143$ \\
MultiFeature & 0.608 & 0.911 & $+50$ \\
\hline
\end{tabularx}
\end{table}

\section*{Заключение}

Проведённое исследование даёт три результата. Долгосрочная память волатильности сильнее и устойчивее, чем в доходностях ($\hat d\approx0.5$ против $\hat d\approx0$). Эта память значимо нестационарна: для 7 из 10 активов выявлены структурные сдвиги, в кризисные периоды сила памяти падает. Локальная оценка памяти несёт самостоятельную прогностическую информацию, улучшая HAR-RV по MSE и QLIKE.

Эти результаты проецируются на устройство трансформерных моделей. Механизм самовнимания играет роль обучаемой автокорреляционной функции, поэтому статичная схема внимания соответствует допущению о постоянной силе памяти. Это допущение нарушается: в кризис память становится короче, и модель, опирающаяся на длинный контекст, ошибается. Отсюда следует модификация внимания~-- сделать его адаптивным к локальной силе памяти: встраивать дифференцируемый оценщик памяти как управляющий сигнал, расширяющий контекст при высокой памяти и сужающий при низкой. Дальнейшая работа предполагает сравнение такой схемы с бейзлайнами и оценку риска VaR/ES.

\ListReferences{
\item Vaswani,~A. Attention is all you need~/ A.~Vaswani [et al.]~// Adv. Neural Inf. Process. Syst.~-- 2017.~-- Vol.~30.~-- P.~5998--6008.
\item Zhou,~H. Informer: Beyond efficient transformer for long sequence time-series forecasting~/ H.~Zhou [et al.]~// Proc. AAAI Conf. on Artificial Intelligence.~-- 2021.~-- Vol.~35, №~12.~-- P.~11106--11115.
\item Hurst,~H.~E. Long-term storage capacity of reservoirs~/ H.~E.~Hurst~// Trans. Am. Soc. Civ. Eng.~-- 1951.~-- Vol.~116, №~1.~-- P.~770--799.
\item Geweke,~J. The estimation and application of long memory time series models~/ J.~Geweke, S.~Porter-Hudak~// J. Time Ser. Anal.~-- 1983.~-- Vol.~4, №~4.~-- P.~221--238.
\item Corsi,~F. A simple approximate long-memory model of realized volatility~/ F.~Corsi~// J. Financial Econometrics.~-- 2009.~-- Vol.~7, №~2.~-- P.~174--196.
}

\end{multicols}

\end{document}
